%  LaTeX support: latex@mdpi.com 
%  For support, please attach all files needed for compiling as well as the log file, and specify your operating system, LaTeX version, and LaTeX editor.

%=================================================================
\documentclass[data,article,submit,pdftex,moreauthors]{Definitions/mdpi} 

%=================================================================
%----------
% journal
%----------

% MDPI internal commands
\firstpage{1} 
\makeatletter 
\setcounter{page}{\@firstpage} 
\makeatother
\pubvolume{1}
\issuenum{1}
\articlenumber{0}
\pubyear{2026}
\copyrightyear{2026}
%\externaleditor{Academic Editor: Firstname Lastname}
\datereceived{06 May 2026} 
\daterevised{20 May 2026} %
\dateaccepted{} 
\datepublished{} 
\hreflink{https://doi.org/} % If needed use \linebreak
\usepackage{tabularx}
\usepackage{booktabs}
\usepackage[flushleft]{threeparttable}
\usepackage{ltablex}
\usepackage{booktabs}
\usepackage{array}
\usepackage{caption}

%=================================================================
\graphicspath{{figures/}}
\usepackage{subcaption}
\usepackage{listings}
\usepackage{xcolor}
\usepackage{textcomp}   % load before gensymb
\usepackage{gensymb}
\usepackage{tabularx}
\newcolumntype{Y}{>{\centering\arraybackslash}X}
\usepackage{amsmath,mathtools,amsthm}
	\setcounter{MaxMatrixCols}{15}
\usepackage{array}
\usepackage{amssymb}
\usepackage{amsfonts}
\usepackage{float}
\usepackage{chngcntr}
\counterwithin*{equation}{section}
\usepackage[super]{nth}
\usepackage{longtable}
\usepackage{tabularx}
\usepackage{multirow}
\usepackage{adjustbox}
\usepackage{soul}  % for \textcolor{blue}{} highlighting

\definecolor{deepblue}{rgb}{0,0,0.5}
\definecolor{deepred}{rgb}{0.6,0,0}
\definecolor{deepmagenta}{RGB}{195,12,214}
\definecolor{deepgreen}{RGB}{29,122,30}
\definecolor{mintgreen}{RGB}{192,249,192}
\definecolor{codepurple}{RGB}{204, 0, 153}
\definecolor{LightCyan}{rgb}{0.88,1,1}
\definecolor{GainsboroPurple}{RGB}{247,176,247}
\definecolor{LightSlateGrey}{RGB}{124,118,118}
%    
\lstdefinestyle{mystyle}{
    commentstyle=\color{LightSlateGrey},
    keywordstyle=\color{deepgreen}\bfseries,
    emphstyle=\color{deepred},
    numberstyle=\tiny\color{deepblue},
    stringstyle=\color{codepurple},
    basicstyle=\ttfamily\scriptsize,
    breakatwhitespace=false,         
    breaklines=true,                 
    captionpos=b,                    
    keepspaces=true,                 
    numbers=left,                    
    numbersep=5pt,                  
    showspaces=false,                
    showstringspaces=false,
    showtabs=false,                  
    tabsize=2
}
\lstset{style=mystyle}


%=================================================================
\Title{Artificial Intelligence and Big Data Analytics for Seismic Hazard Assessment: \textcolor{blue}{Methodological Advances and Computational Frameworks for the} Marmara Region, Türkiye}

\TitleCitation{Artificial Intelligence and Big Data Analytics for Seismic Hazard Assessment: \textcolor{blue}{Methodological Advances and Computational Frameworks} for the Marmara Region, Türkiye}

\newcommand{\orcidauthorA}{0000-0002-5759-1089}

\Author{Polina Lemenkova $^{1}$*\orcidA{}, Abdullah Can Zülfikar $^{1}$}

\AuthorNames{Polina Lemenkova, Abdullah Can Zülfikar}

\AuthorCitation{Lemenkova, P.; Zülfikar, A.C.} 

\address{%
$^{1}$ \quad Disaster and Emergency Management Department, Institute of Earthquake Engineering and Disaster Management, Istanbul Technical University. İTÜ Ayazağa, 34469, Maslak Kampüsü, İstanbul, Türkiye.
}

\corres{Correspondence: lemenkova@itu.edu.tr}

\abstract{\textcolor{blue}{The Marmara region of Türkiye, situated along the North Anatolian Fault Zone (NAFZ), constitutes one of the most seismically active and densely monitored zones globally.} Given the region's high vulnerability and the catastrophic impacts of historical events—\textcolor{blue}{notably} the 1999 İzmit and 2023 Kahramanmaraş sequences—there is a critical need for advanced seismic hazard \textcolor{blue}{risk} assessment \textcolor{blue}{(SHRA) methods that move beyond static models.} This \textcolor{blue}{review} \textcolor{blue}{examines} the paradigm shift from traditional geophysics to big data seismology, \textcolor{blue}{characterized by the ``Five Vs'': volume, velocity, variety, veracity, and value.} \textcolor{blue}{Critically, we distinguish between two fundamentally different problems: Earthquake Early Warning (EEW), which operates on sub-second timescales after rupture initiation, and probabilistic earthquake forecasting, which operates on timescales of years to decades.} The study \textcolor{blue}{discusses} how \textcolor{blue}{cloud-native platforms such as} Azure Databricks\textcolor{blue}{, combined with} data pipelines using Apache Kafka and Apache Spark\textcolor{blue}{, enable the} real-time processing of petabyte-scale \textcolor{blue}{seismic} sensor streams\textcolor{blue}{.} Key technological \textcolor{blue}{tools}, including Physics-Informed Neural Networks (PINNs) and \textcolor{blue}{deep learning} models \textcolor{blue}{such as} PhaseNet, are analyzed for their \textcolor{blue}{demonstrated} ability to enhance \textcolor{blue}{EEW} systems \textcolor{blue}{through} sub-second phase picking and automated event detection. \textcolor{blue}{We present statistical validation metrics and uncertainty quantification methods essential for credible hazard assessment.} By addressing computational bottlenecks through hybrid \textcolor{blue}{computing} architectures and edge computing, this framework aims to \textcolor{blue}{improve the warning lead time for} Istanbul's critical infrastructure. \textcolor{blue}{This work provides a structured roadmap} for bridging the gap between traditional \textcolor{blue}{seismic} data analysis and \textcolor{blue}{operational} predictive analytics \textcolor{blue}{in the Marmara region.}}

\keyword{big data; seismology; cloud computing; geophysics; \textcolor{blue}{machine learning}; \textcolor{blue}{earthquake early warning}; signal processing; \textcolor{blue}{Apache Spark; North Anatolian Fault}} 

\PACS{91.30.Ab, 91.30.Pd, 07.05.Mh, 89.20.Ff}
\MSC{86A15, 68T07}
\JEL{Q54, C45, O33}

\begin{document}

%----------- section ----------->
\section{Introduction}

%----------- section ----------->
\subsection{Background}

Seismic hazard monitoring systems prioritize the early detection \textcolor{blue}{and} probabilistic assessment of earthquake events to mitigate geotechnical risks \cite{Allen2012,Kanamori2005}. Modern geophysical centers analyze thousands of global sensor streams per second, generating millions of daily data points \cite{USGS2020,IRIS2018,Chung2020}. \textcolor{blue}{However, the transition to operational seismic monitoring necessitates} processing capabilities that exceed the limits of traditional geophysical analysis \cite{Beroza2018,Allen2019}. The integration of Artificial Intelligence (AI) and cloud computing \textcolor{blue}{has fundamentally} transformed this landscape, enabling the treatment of global seismic signals as \textcolor{blue}{unified} big data \textcolor{blue}{repositories} rather than isolated regional records \cite{Perol2018,Zhu2019,Kong2019}.

The paradigm shift from traditional geophysics to big data seismology \textcolor{blue}{represents} a fundamental change in analytical methodology \cite{Bergen2019,Donner2019}. Contemporary seismic datasets require novel frameworks for storage and high-throughput processing to extract \textcolor{blue}{reliable} prognostic insights \cite{Ahern2018,IRISDMC2015}. \textcolor{blue}{Crucially, modern seismic networks demand a departure from legacy data management systems toward} distributed, scalable data science architectures, Figure~\ref{fig01}.

\begin{figure}[H]
    \begin{center}
	\hspace{-35pt}\includegraphics[width=13.0cm]{Fig_01.png}
    \end{center}
    \caption{Big Data in Geophysics and Seismology: the 5Vs framework. Diagram source: authors.}
    \label{fig01}
\end{figure}

The ``Five Vs'' of big data—volume, velocity, variety, veracity, and value—define the \textcolor{blue}{principal} advantages over regional geophysical databases \cite{Laney2001,Kitchin2014,Hashem2015}. While traditional archives typically manage gigabytes of localized \textcolor{blue}{data}, modern big data systems routinely ingest petabytes of global sensor data \cite{Ahern2018}. \textcolor{blue}{The} velocity of these systems enables the processing of real-time data streams, providing \textcolor{blue}{the} immediate \textcolor{blue}{and} latency-critical insights required for seismic risk evaluation and rapid decision-making \cite{gandomi2015beyond,hashem2015rise,allen2009earthquake}.

The variety of big data in geophysics is \textcolor{blue}{defined} by a spectrum of heterogeneous data structures \cite{li2016big,wang2019geophysical}. The complexity of formats arises \textcolor{blue}{from} the diverse data structures and sources \textcolor{blue}{that} must be integrated and \textcolor{blue}{jointly} analysed \cite{chen2014big}. \textcolor{blue}{Even tabular CSV data typically wraps heterogeneous information pulled} from \textcolor{blue}{multiple} geophysical sensors \cite{ma2015big,kessler2019data}. \textcolor{blue}{Finally,} the veracity \textcolor{blue}{dimension} \textcolor{blue}{pertains} to data quality and integrity \cite{fang2015big,reyes2017data}. \textcolor{blue}{Seismic datasets} may include incomplete signals, duplicate \textcolor{blue}{event} entries, inconsistent formatting, \textcolor{blue}{malfunctioning} sensors, or outdated information \cite{mousavi2019cred}.

\textcolor{blue}{These five dimensions create interconnected challenges that conventional seismological methods are ill-equipped to handle.} Traditional analysis \textcolor{blue}{largely} relies on manual examination of datasets, \textcolor{blue}{neglecting} massive volumes of historical seismic data \textcolor{blue}{available in} archives and \textcolor{blue}{from} sensors \textcolor{blue}{flowing into} online geophysical repositories \cite{kong2019machine}. \textcolor{blue}{Processing such large data massifs enables the discovery of} anomalies in mantle processes with better precision and speed\textcolor{blue}{, yielding} real-time data characterizing location and magnitude of seismic events, \textcolor{blue}{together with} unstructured cartographic data and \textcolor{blue}{relevant} geological characteristics \cite{perol2018convolutional,zhu2019phasenet}. \textcolor{blue}{Consequently,} big data analytics in seismology \textcolor{blue}{support} predictive analytics that \textcolor{blue}{process} sensor data from thousands of \textcolor{blue}{stations} simultaneously.

%----------- section ----------->
\subsection{Research Problem and Objectives}

\textcolor{blue}{Big data technologies are revolutionizing geophysics by enabling dynamic, high-resolution risk assessment.} In Seismic Hazard Risk Assessment (SHRA), big data \textcolor{blue}{has shifted the paradigm} from static, \textcolor{blue}{low-frequency} maps to dynamic\textcolor{blue}{,} real-time monitoring models. Data integration and real-time processing \textcolor{blue}{allow} historical data analysis, current data monitoring\textcolor{blue}{,} and future data forecasting \cite{Sen2009}. Recent decades have witnessed diverse geological \textcolor{blue}{and} climate hazards \textcolor{blue}{whose} cumulative effects on the environment \textcolor{blue}{are well-documented} \cite{Liggins2013}. \textcolor{blue}{Seismic risks are particularly prominent} due to their devastating consequences for society, \textcolor{blue}{demanding} effective \textcolor{blue}{quantitative} seismic risk assessment methods \cite{Dowrick2009,Liu2024}.

\textcolor{blue}{Before proceeding, it is essential to distinguish between two fundamentally different scientific problems that are often conflated in the literature: (i) Earthquake Early Warning (EEW) operates on timescales of seconds, detecting P-waves from an initiated rupture and issuing alerts before the destructive S-waves arrive; and (ii) earthquake forecasting operates on timescales of years to decades, using probabilistic seismicity models to estimate long-term rupture probability. These two paradigms require distinct data architectures, latency constraints, and validation frameworks.}

This paper \textcolor{blue}{contributes} to the challenge of big data technologies in geophysics \textcolor{blue}{by analyzing} how data \textcolor{blue}{pipelines} \textcolor{blue}{support} modern \textcolor{blue}{operational} seismic hazard risk assessment. \textcolor{blue}{We discuss SHRA in the context of big data and high-scale geophysical analytics, bridging} the gap between traditional data analysis in seismology and predictive analytics. \textcolor{blue}{Unlike conventional approaches that rely on processing single CSV files, a modern} big data \textcolor{blue}{pipeline} can \textcolor{blue}{architect} an SHRA system that handles \textcolor{blue}{the} volume\textcolor{blue}{s}, velocity\textcolor{blue}{,} and variety of seismic signals simultaneously.

%----------- section ----------->

\section{Seismic Vulnerability in Türkiye: From Tectonic Plate Interactions to Historical Disaster Impacts}

\subsection{Geology and Tectonics in Türkiye}

Türkiye is \textcolor{blue}{highly} vulnerable to seismic \textcolor{blue}{hazards}, especially in regions near active tectonic faults \cite{Basaran2014,Kandregula2024}. \textcolor{blue}{Systematic} geological monitoring is essential for comprehensive seismic risk assessment, underscoring the need for \textcolor{blue}{improved} methodologies \cite{Korkmaz2013}. \textcolor{blue}{Seismically active areas in Türkiye lie at the margins of the Anatolian Block} \cite{Jolivet2023}, within a zone of lithosphere collision \textcolor{blue}{involving} three major tectonic plates: the Eurasian Plate, the African Plate, and the Arabian Plate.

\textcolor{blue}{The Anatolian Block is driven westward as the Arabian Plate advances northward against the Eurasian Plate, while the African Plate subducts beneath the Hellenic Arc} \cite{McKenzie1972,LePichon2010}. \textcolor{blue}{These plate interactions generate} frequent and powerful earthquakes \textcolor{blue}{at tectonic boundaries}. The most prominent feature is the North Anatolian Fault Zone (NAFZ), a 1\,200~km right-lateral strike-slip fault extending from the Gulf of Saros to Karlıova \cite{Barka1992,Sengor2005}. \textcolor{blue}{Formed 13--11 million years ago}, it remains seismically active\textcolor{blue}{,} capable of producing earthquakes \textcolor{blue}{up} to magnitude 8.0 \cite{Stein1997}. Current deformation, particularly in the Marmara region, is asymmetric and \textcolor{blue}{strongly} influenced by \textcolor{blue}{the} complex regional geology \textcolor{blue}{of the NAFZ} \cite{Puccinelli2021}.

\subsection{Historical Outline of Seismic Events in Türkiye}

The historical record of the Anatolian region reveals a relentless cycle of seismic activity, Figure~\ref{fig03}.

\begin{figure}[H]
    \begin{center}
	\hspace{-35pt}\includegraphics[width=14.0cm]{Fig_03.png}
    \end{center}
    \caption{Historical outline of \textcolor{blue}{major} seismic events in Türkiye. Diagram source: authors.}
    \label{fig03}
\end{figure}

In the Aegean region, from 496~BC to 1949~AD, the area experienced a total of 20 moderate to severe earthquakes \cite{altinok2011,papazachos2000}. Many ancient events resulted in devastating tsunamis, with notable occurrences in 1389, 1856, 1866, 1881, and 1949 \cite{yalciner2002}. The 1881 Chios-\c{C}e\c{s}me earthquake \textcolor{blue}{caused} thousands of fatalities and significant coastal changes \cite{ambraseys2000}. The $M_w$~7.4 İzmit earthquake of 17 August 1999 \textcolor{blue}{remains one of the best-studied large events in the region} \cite{Altinoketal2001,Inanetal2007}. \textcolor{blue}{Parsons (2004) estimated a} significantly heightened probability (\textcolor{blue}{$\sim$}50\%) of \textcolor{blue}{a} $M \geq 7$ event \textcolor{blue}{within} 30 \textcolor{blue}{years} \textcolor{blue}{of the 1999 stress transfer} \cite{Parsons2004}.

Most recently, on February~6, 2023, Türkiye experienced a catastrophic seismic sequence with a mainshock of $M_w$~7.8, followed by \textcolor{blue}{major} earthquakes of $M_w$~7.7 and 7.6. \textcolor{blue}{These events caused total destruction in Antakya} \cite{Akcan2024,Akinci2025}.

\subsection{Seismic Risk Assessment and the Distinction Between EEW and Earthquake Forecasting}

\textcolor{blue}{The 2023 Kahramanmaraş sequence reinforced the urgent need to clearly separate two domains of operational seismology that serve different societal functions.} Research and technological development in \textcolor{blue}{EEW systems have become critical components of seismic risk mitigation} in the Anatolian region. \textcolor{blue}{EEW is fundamentally a real-time detection and alerting problem: once a rupture has begun, P-wave signals are detected, source parameters are estimated within seconds, and automated alerts are issued before destructive shaking arrives. In contrast, earthquake forecasting—such as the Parsons (2004) probability estimate—is a probabilistic, long-term assessment of where and when future earthquakes are likely to occur, based on fault mechanics and historical seismicity.}

\textcolor{blue}{This distinction matters for data infrastructure: EEW requires deterministic, sub-second pipelines with strict latency constraints ($<$100~ms), while forecasting models require access to decades of historical catalogs, geodetic strain data, and paleoseismic records.} The Kandilli Observatory and Earthquake Research Institute (KOERI) \textcolor{blue}{has} operationalized a new-generation EEW\textcolor{blue}{S} across Türkiye\textcolor{blue}{,} leveraging a dense network of stations to issue alerts within seconds \cite{Ozeletal2026}. Innovative strategies \textcolor{blue}{also} explore hybrid models, such as using mosque loudspeakers to disseminate warnings \cite{Eren2025}. \textcolor{blue}{Machine learning is} transforming \textcolor{blue}{raw telemetry} into actionable protocols \cite{Doganetal2019}, \textcolor{blue}{predicting} potential damage zones and automatically triggering safety actions \cite{CremenGalasso2020,Specialeetal2026}.

\section{Research Gap Identification and Contribution}

\textcolor{blue}{To our knowledge, no} comprehensive review \textcolor{blue}{exists that integrates big data processing architectures, statistical validation frameworks, and the EEW/forecasting distinction specifically for the Marmara region.} \textcolor{blue}{While individual} case studies \textcolor{blue}{on AI in seismology exist} \cite{Bergen2019,MousaviBeroza2022}, \textcolor{blue}{integration reviews addressing big data pipelines, uncertainty quantification, and cloud governance for Turkish seismic monitoring are lacking.} This \textcolor{blue}{article} provides a comprehensive analysis of the evolution of techniques for big data \textcolor{blue}{processing} in the context of seismic risk assessment of Türkiye, \textcolor{blue}{with a focus on the Marmara region and the NAFZ.} Mitigating earthquake risks \textcolor{blue}{requires both} an EEW system \textcolor{blue}{and} consistent \textcolor{blue}{probabilistic} seismic monitoring \textcolor{blue}{underpinned by} big data, Figure~\ref{fig04}.

\begin{figure}[H]
    \begin{center}
	\hspace{-35pt}\includegraphics[width=14.0cm]{Fig_04.png}
    \end{center}
    \caption{Development of Big Data in Seismology, EEW \& Geophysical Monitoring in Istanbul, Türkiye. Diagram source: authors.}
    \label{fig04}
\end{figure}

\textcolor{blue}{The current limitations of conventional seismological processing are summarized in Table~\ref{tab01}, which highlights both the technical constraints and their implications for Marmara region monitoring.}

\begin{table}[ht]
\centering
\begin{threeparttable}
\caption{\textcolor{blue}{Limitations of Existing Geophysical Processing Approaches and Their Marmara-Region Implications}}
\label{tab01}
\begin{tabularx}{\textwidth}{@{} l l X @{}}
\toprule
\textbf{Approach} & \textbf{Core Constraint} & \textbf{Big Data \& Geophysical Impact} \\ \midrule
Traditional Inversion & Non-uniqueness & Multiple subsurface models satisfy the same observed data. \textcolor{blue}{High-dimensionality exacerbates uncertainty quantification; confidence intervals on hazard maps are rarely reported.} \\ \addlinespace
Linearized Models & Oversimplification & Standard kernels assume a linear Earth. \textcolor{blue}{Complex NAFZ geology is highly non-linear; forcing linearity produces false anomalies and signal artifacts in high-volume datasets.} \\ \addlinespace
Manual QC & Human Bottleneck & \textcolor{blue}{High-density 3D/4D seismic surveys generate petabytes of traces.} Manual QC is the primary latency in the processing pipeline \textcolor{blue}{and is infeasible for continuous monitoring.} \\ \addlinespace
HPC Constraints & Data Movement & In seismic migration (e.g., RTM), moving massive velocity models from storage to compute nodes consumes more energy than the FLOPs performed (Von Neumann bottleneck\tnote{1}). \\ \bottomrule
\end{tabularx}
\begin{tablenotes}
\small
\item[1] The Von Neumann bottleneck refers to the limited throughput between the central processing unit (CPU) and memory compared to the amount of memory.
\end{tablenotes}
\end{threeparttable}
\end{table}

%----------- section ----------->
\section{Big Data Processing in Geophysics}

%----------- subsection ----------->
\subsection{Hybrid Computing: HPC and Cloud Architecture}

One of the most demanding data processing tasks in geophysics involves massive datasets from nodal sensors and complex inversions, such as Full Waveform Inversion (FWI) \cite{witte2019scalable}. Since nodal datasets can reach petabyte scales, geophysicists \textcolor{blue}{no longer move all data to a single location; instead,} they use a hybrid \textcolor{blue}{computing} approach \cite{peter2011forward,wang2020cloud}, \textcolor{blue}{combining} the advantages of both HPC and cloud computing \cite{fox2017big,assuncao2015big}.

Handling big data in geophysics requires a shift from \textcolor{blue}{sequential} processing to a hybrid architecture that blends high-performance computing (HPC) with cloud-native intelligence \cite{fox2017big,hashem2015rise}. To process massive nodal datasets\textcolor{blue}{,} geophysicists utilise a hybrid architecture combining GPU-accelerated HPC clusters with elastic cloud scalability, \textcolor{blue}{moving the processing code to the data to avoid massive transfer bottlenecks} \cite{wang2020cloud}. FWI \textcolor{blue}{remains on GPU-accelerated HPC clusters due to the tight coupling required between processors. Bursty tasks such as initial signal denoising are offloaded to elastic cloud resources.} The hybrid infrastructure accelerates high-resolution FWI by embedding physical laws into \textcolor{blue}{Physics-Informed Neural Networks (PINNs)} \cite{raissi2019physics,sun2020surrogate}.

\textcolor{blue}{To quantify the performance of this hybrid approach, we adopt the following statistical validation framework for Marmara-region EEW pipelines:}

\begin{equation}
\text{Precision} = \frac{TP}{TP + FP}, \quad \text{Recall} = \frac{TP}{TP + FN}
\end{equation}

\begin{equation}
F_1 = 2 \cdot \frac{\text{Precision} \times \text{Recall}}{\text{Precision} + \text{Recall}}
\end{equation}

\textcolor{blue}{where $TP$, $FP$, and $FN$ denote true positives, false positives, and false negatives in event detection, respectively. For EEW, false negatives (missed alerts) carry higher societal cost than false positives (false alarms).}

\begin{equation}
\textcolor{blue}{\text{RMSE}(\hat{M}, M) = \sqrt{\frac{1}{N}\sum_{i=1}^{N}(\hat{M}_i - M_i)^2}}
\end{equation}

\textcolor{blue}{This root-mean-square error metric is used to evaluate magnitude estimation accuracy across $N$ detected events, where $\hat{M}_i$ and $M_i$ denote estimated and catalog magnitudes, respectively.}

\textcolor{blue}{Uncertainty in the estimated epicentral distance $\Delta r$ is propagated through the EEW warning time estimate:}
\begin{equation}
\textcolor{blue}{T_{\text{warn}} = \frac{\Delta r}{V_S} - T_{\text{proc}}}
\end{equation}

\textcolor{blue}{where $V_S$ is the S-wave velocity and $T_{\text{proc}}$ is the end-to-end processing latency. For Istanbul, with $\Delta r \approx 50$~km from the NAFZ Marmara segment and $V_S \approx 3.5$~km/s, $T_{\text{warn}} \leq 14.3$~s under ideal conditions.}

For remote nodal arrays, smart nodes perform initial data compression and basic quality checks at the source, Table~\ref{tab02}.

\keepXColumns 

\begin{tabularx}{\textwidth}{@{} >{\hsize=0.8\hsize\raggedright\arraybackslash}X >{\hsize=1.2\hsize\raggedright\arraybackslash}X >{\hsize=1.0\hsize\raggedright\arraybackslash}X @{}}
    \caption{State-of-the-Art Big Data Processing in Geophysics \textcolor{blue}{(2026 Framework)}} \label{tab02} \\
    \toprule
    \textbf{Geophysical Task} & \textbf{\textcolor{blue}{Current} Advanced Strategy} & \textbf{Technology Stack} \\ 
    \midrule
    \endfirsthead   
    \multicolumn{3}{c}{{\bfseries \tablename\ \thetable{} -- Continued from previous page}} \\
    \toprule
    \textbf{Geophysical Task} & \textbf{\textcolor{blue}{Current} Advanced Strategy} & \textbf{Technology Stack} \\ 
    \midrule
    \endhead    
    \midrule
    \multicolumn{3}{r}{{Continued on next page...}} \\
    \endfoot    
    \bottomrule
    \endlastfoot
    \textbf{High-Density Nodal Ingest} & Decoupled asynchronous streaming from 100k+ channel nodes; real-time metadata indexing. & Azure Blob/S3, Apache Kafka, Zarr/ASDF Formats \\ \addlinespace
    
    \textbf{Elastic FWI \& Velocity Modeling} & Multi-parameter PINN-based inversion to resolve complex salt geometries and anisotropy. & NVIDIA H200 Clusters, PyTorch/JAX-Seismic \\ \addlinespace
    
    \textbf{Seismic Signal Enhancement} & Self-supervised Deep Learning Autoencoders for 5D interpolation and ghost reflection removal. & Transformer-based Denoising, GANs \\ \addlinespace
    
    \textbf{Microseismic Monitoring} & Automated event detection and location using Edge-AI for real-time hydraulic fracture mapping. & NVIDIA Jetson AGX Orin, 5G Telemetry \\ \addlinespace
    
    \textbf{4D Reservoir Characterization} & Digital Twin synchronization integrating seismic, EM, and production data for predictive flow modeling. & Azure Digital Twins, NVIDIA Omniverse \\ \addlinespace
    
    \textbf{Multi-Physics Integration} & Joint inversion of gravity, magnetic, and seismic data via \textcolor{blue}{foundation} models for \textcolor{blue}{geophysical} interpretation. & Foundation Models, Graph Neural Networks (GNNs) \\ 
\end{tabularx}

%----------- subsection ----------->
\subsection{Physics-Informed Neural Networks (PINNs) and Deep Learning for Phase Picking}

The most demanding inversions are now handled by PINNs \cite{Aghamiry2021}. Unlike traditional AI, \textcolor{blue}{which} might suggest geologically impossible structures, PINNs embed the actual wave equations (the physics) into the neural network's loss function. \textcolor{blue}{The governing acoustic wave equation embedded in the PINN loss is:}

\begin{equation}
\textcolor{blue}{\frac{\partial^2 u}{\partial t^2} - v^2(\mathbf{x})\nabla^2 u = f(\mathbf{x}, t)}
\end{equation}

\textcolor{blue}{where $u(\mathbf{x},t)$ is the wavefield, $v(\mathbf{x})$ is the P-wave velocity model, and $f(\mathbf{x},t)$ is the seismic source function. The total PINN loss function combines data misfit and physics residual:}

\begin{equation}
\textcolor{blue}{\mathcal{L}_{\text{total}} = \mathcal{L}_{\text{data}} + \lambda \mathcal{L}_{\text{physics}}}
\end{equation}

\textcolor{blue}{where $\lambda$ is a weighting hyperparameter controlling the trade-off between data fitting and physical consistency.}

\textcolor{blue}{For phase picking, the convolutional neural network PhaseNet} \cite{ZhuBeroza2019,Park2020,Ni2023} \textcolor{blue}{achieves a reported P-wave pick precision exceeding 95\% on the STEAD benchmark dataset, reducing manual analyst time by several orders of magnitude.}

\textcolor{blue}{The uncertainty in P-wave arrival time $\sigma_t$ propagates directly into epicentral location uncertainty:}

\begin{equation}
\textcolor{blue}{\sigma_r = V_P \cdot \sigma_t}
\end{equation}

\textcolor{blue}{For $\sigma_t = 0.05$~s (typical for PhaseNet) and $V_P = 6.0$~km/s, $\sigma_r \approx 0.3$~km, which is adequate for EEW source parameter estimation.}

%----------- section ----------->
\section{Cluster Computing and Distributed Data Processing}

%----------- subsection ----------->
\subsection{Architecture Paradigms for Marmara Regional Seismology}

To manage the heterogeneous data demands of the Marmara region, cluster architectures are categorized into three patterns. The Master-Worker paradigm, utilizing Apache Spark, facilitates \textcolor{blue}{high-throughput} batch analytics essential for computationally intensive FWI and 3D tomographic modeling of the NAFZ. For high-density nodal arrays, a Peer-to-Peer pattern (e.g., Apache Cassandra) employs decentralized gossip protocols and N-replication to ensure fault tolerance. \textcolor{blue}{A} Hybrid Broker-Based architecture using Apache Kafka provides the decoupled scaling necessary to ingest \textcolor{blue}{and} partition massive telemetry streams, Figure~\ref{fig04}.

%----------- subsection ----------->
\subsection{Apache Spark and PySpark for Seismic Data Processing}

\textcolor{blue}{The following eight code scripts illustrate the practical application of Apache Spark and related Python tools in big data seismology. These examples address the full pipeline from raw ingestion to validated event detection.}

\textbf{Script~1: Reading and Windowing a Seismic Waveform Stream with PySpark}

\begin{lstlisting}[language=Python, caption={PySpark ingestion and windowing of continuous seismic waveform data from AFAD-compatible MiniSEED streams.}]
from pyspark.sql import SparkSession
from pyspark.sql.functions import window, col, from_json
from pyspark.sql.types import StructType, StringType, DoubleType, TimestampType

spark = SparkSession.builder \
    .appName("MarmaraSeismicIngest") \
    .config("spark.sql.shuffle.partitions", "200") \
    .getOrCreate()

schema = StructType() \
    .add("station_id", StringType()) \
    .add("timestamp", TimestampType()) \
    .add("channel", StringType()) \
    .add("amplitude", DoubleType()) \
    .add("sampling_rate", DoubleType())

# Read streaming data from Apache Kafka (Bronze Layer)
raw_stream = spark.readStream \
    .format("kafka") \
    .option("kafka.bootstrap.servers", "kafka-broker:9092") \
    .option("subscribe", "marmara.seismic.raw") \
    .load() \
    .selectExpr("CAST(value AS STRING) as json_str") \
    .select(from_json(col("json_str"), schema).alias("data")) \
    .select("data.*")

# Apply 10-second tumbling windows for real-time analysis
windowed = raw_stream \
    .groupBy(window(col("timestamp"), "10 seconds"), col("station_id")) \
    .agg({"amplitude": "max"}) \
    .withColumnRenamed("max(amplitude)", "peak_amplitude")

query = windowed.writeStream \
    .outputMode("append") \
    .format("delta") \
    .option("checkpointLocation", "/mnt/checkpoints/seismic_ingest") \
    .start("/mnt/bronze/seismic_windows")
query.awaitTermination()
\end{lstlisting}

\textbf{Script~2: Butterworth Bandpass Filtering in the Silver Layer}

\begin{lstlisting}[language=Python, caption={Application of a Butterworth bandpass filter to seismic waveforms using ObsPy within a PySpark UDF to remove cultural noise.}]
import numpy as np
from scipy.signal import butter, sosfilt
from pyspark.sql.functions import udf
from pyspark.sql.types import ArrayType, DoubleType

def butterworth_bandpass(signal_list, lowcut=1.0, highcut=10.0, fs=100.0, order=4):
    """Apply a Butterworth bandpass filter to isolate tectonic signal."""
    sos = butter(order, [lowcut, highcut], btype='band', fs=fs, output='sos')
    return sosfilt(sos, np.array(signal_list)).tolist()

bandpass_udf = udf(butterworth_bandpass, ArrayType(DoubleType()))

silver_df = spark.read.format("delta").load("/mnt/bronze/seismic_windows")
filtered_df = silver_df.withColumn(
    "filtered_signal",
    bandpass_udf(col("raw_samples"), col("sampling_rate"))
)
filtered_df.write.format("delta") \
    .mode("overwrite") \
    .save("/mnt/silver/filtered_waveforms")
print(f"Processed {filtered_df.count()} waveform records into Silver layer.")
\end{lstlisting}

\textbf{Script~3: STA/LTA Trigger-Based Event Detection}

\begin{lstlisting}[language=Python, caption={Short-Term Average / Long-Term Average (STA/LTA) event detection implemented as a distributed PySpark operation across the AFAD station network.}]
from pyspark.sql.functions import pandas_udf
from pyspark.sql.types import BooleanType
import pandas as pd

@pandas_udf(BooleanType())
def sta_lta_trigger(signal_series: pd.Series) -> pd.Series:
    """
    Detect seismic events using the STA/LTA ratio.
    Returns True where STA/LTA exceeds threshold (typically 3.0).
    """
    results = []
    for signal in signal_series:
        arr = np.array(signal)
        sta_len, lta_len, threshold = 5, 100, 3.0
        sta = np.convolve(arr**2, np.ones(sta_len)/sta_len, 'same')
        lta = np.convolve(arr**2, np.ones(lta_len)/lta_len, 'same')
        ratio = np.where(lta > 0, sta / lta, 0)
        results.append(bool(np.max(ratio) > threshold))
    return pd.Series(results)

events_df = filtered_df.withColumn(
    "event_detected", sta_lta_trigger(col("filtered_signal"))
).filter(col("event_detected") == True)

print(f"Detected {events_df.count()} candidate seismic events.")
events_df.write.format("delta").mode("append") \
    .save("/mnt/silver/detected_events")
\end{lstlisting}

\textbf{Script~4: Magnitude Estimation and Statistical Validation}

\begin{lstlisting}[language=Python, caption={Local magnitude estimation using peak amplitude and statistical validation (RMSE, precision, recall) against the AFAD earthquake catalog.}]
import math
from pyspark.sql.functions import log10, sqrt, mean, stddev

# Local magnitude: Ml = log10(A) + 1.11*log10(R) + 0.00189*R - 2.09
# where A = peak amplitude (nm), R = hypocentral distance (km)
catalog_df = spark.read.format("delta").load("/mnt/gold/afad_catalog")
predicted_df = events_df \
    .withColumn("log_amp", log10(col("peak_amplitude"))) \
    .withColumn("Ml_estimated",
        col("log_amp") + 1.11 * log10(col("distance_km"))
        + 0.00189 * col("distance_km") - 2.09)

# Join with catalog for validation
joined = predicted_df.join(catalog_df, on="event_id", how="inner")
validation = joined.withColumn("sq_error",
    (col("Ml_estimated") - col("Ml_catalog"))**2)

rmse = validation.agg(sqrt(mean(col("sq_error"))).alias("RMSE")).collect()[0]["RMSE"]
print(f"Magnitude estimation RMSE: {rmse:.3f}")

# Precision / Recall for event detection (TP threshold: |dM| < 0.5)
tp = joined.filter((col("Ml_estimated") - col("Ml_catalog")).between(-0.5, 0.5)).count()
fp = predicted_df.count() - tp
fn = catalog_df.count() - tp
precision = tp / (tp + fp) if (tp + fp) > 0 else 0
recall    = tp / (tp + fn) if (tp + fn) > 0 else 0
f1        = 2 * precision * recall / (precision + recall) if (precision + recall) > 0 else 0
print(f"Precision={precision:.3f}, Recall={recall:.3f}, F1={f1:.3f}")
\end{lstlisting}

\textbf{Script~5: Ambient Noise Cross-Correlation for Surface Wave Tomography}

\begin{lstlisting}[language=Python, caption={Distributed ambient noise cross-correlation for surface wave tomography using PySpark's cartesian join across all Marmara station pairs.}]
from pyspark.sql.functions import udf
from pyspark.sql.types import ArrayType, DoubleType
import numpy as np

@udf(returnType=ArrayType(DoubleType()))
def cross_correlate(sig1, sig2):
    """Normalized cross-correlation for ambient noise interferometry."""
    a, b = np.array(sig1), np.array(sig2)
    a = a / (np.std(a) + 1e-10)
    b = b / (np.std(b) + 1e-10)
    xcorr = np.correlate(a, b, mode='full')
    return (xcorr / len(a)).tolist()

stations_df = spark.read.format("delta").load("/mnt/silver/daily_noise")
pairs_df = stations_df.alias("s1").crossJoin(stations_df.alias("s2")) \
    .filter(col("s1.station_id") < col("s2.station_id"))

xcorr_df = pairs_df.withColumn(
    "xcorr",
    cross_correlate(col("s1.waveform"), col("s2.waveform"))
)
xcorr_df.write.format("delta").mode("append") \
    .save("/mnt/gold/noise_xcorr_marmara")
print("Cross-correlation complete for all Marmara station pairs.")
\end{lstlisting}

\textbf{Script~6: PhaseNet-Based Automated P/S Phase Picking with Uncertainty}

\begin{lstlisting}[language=Python, caption={Automated P- and S-wave phase picking using PhaseNet integrated into a PySpark Gold-layer pipeline with pick uncertainty estimation.}]
import torch
import numpy as np
from pyspark.sql.functions import struct, col

# Load pre-trained PhaseNet model (seisbench interface)
# seisbench.models.PhaseNet.from_pretrained("stead")
def phasenet_pick(waveform_list):
    """
    Apply PhaseNet to 3-component waveform and return
    (p_pick_sample, s_pick_sample, p_prob, s_prob).
    """
    wf = np.array(waveform_list).reshape(3, -1)
    # Normalize per channel
    wf = (wf - wf.mean(axis=1, keepdims=True)) / (wf.std(axis=1, keepdims=True) + 1e-8)
    with torch.no_grad():
        probs = model(torch.tensor(wf, dtype=torch.float32).unsqueeze(0))
    p_prob = probs[0, 1, :].numpy()
    s_prob = probs[0, 2, :].numpy()
    p_pick = int(np.argmax(p_prob))
    s_pick = int(np.argmax(s_prob))
    return (p_pick, s_pick, float(p_prob[p_pick]), float(s_prob[s_pick]))

schema_pick = "struct<p_pick:int, s_pick:int, p_prob:double, s_prob:double>"
phasenet_udf = udf(phasenet_pick, schema_pick)

picked_df = filtered_df \
    .withColumn("picks", phasenet_udf(col("filtered_signal"))) \
    .select("station_id", "timestamp",
            col("picks.p_pick").alias("p_sample"),
            col("picks.s_pick").alias("s_sample"),
            col("picks.p_prob").alias("p_confidence"),
            col("picks.s_prob").alias("s_confidence"))

# Uncertainty: sigma_t = (1 - confidence) * dt, dt = 1/sampling_rate
picked_df = picked_df.withColumn(
    "p_uncertainty_s", (1.0 - col("p_confidence")) / 100.0)
picked_df.write.format("delta").mode("append") \
    .save("/mnt/gold/phase_picks")
\end{lstlisting}

\textbf{Script~7: Real-Time EEW Alert Pipeline with Latency Monitoring}

\begin{lstlisting}[language=Python, caption={End-to-end EEW alert pipeline using Spark Structured Streaming, monitoring processing latency to maintain the ``Golden Seconds'' constraint.}]
from pyspark.sql.functions import current_timestamp, unix_timestamp, expr

alert_stream = picked_df \
    .filter(col("p_confidence") > 0.85) \
    .withColumn("ingestion_time", current_timestamp()) \
    .withColumn("latency_ms",
        (unix_timestamp(col("ingestion_time")) -
         unix_timestamp(col("timestamp"))) * 1000)

# Flag records exceeding 100 ms EEW latency budget
alert_stream = alert_stream \
    .withColumn("latency_ok", col("latency_ms") < 100)

# Write alerts to dashboard and trigger automated responses
alert_query = alert_stream.writeStream \
    .outputMode("append") \
    .format("delta") \
    .option("checkpointLocation", "/mnt/checkpoints/eew_alerts") \
    .trigger(processingTime="1 second") \
    .start("/mnt/gold/eew_alerts")

# Log latency SLA compliance
sla_df = spark.read.format("delta").load("/mnt/gold/eew_alerts")
sla_pct = sla_df.filter(col("latency_ok")).count() / sla_df.count() * 100
print(f"EEW latency SLA compliance: {sla_pct:.1f}%")
alert_query.awaitTermination()
\end{lstlisting}

\textbf{Script~8: Probabilistic Seismic Hazard Estimation Using Gutenberg--Richter Statistics}

\begin{lstlisting}[language=Python, caption={Gutenberg--Richter b-value estimation and exceedance probability computation from the AFAD catalog using PySpark SQL and NumPy.}]
import numpy as np
from pyspark.sql.functions import count, col, log10 as spark_log10

# Gutenberg-Richter: log10(N) = a - b*M
catalog = spark.read.format("delta").load("/mnt/gold/afad_catalog")

mag_counts = catalog \
    .filter(col("magnitude") >= 2.0) \
    .groupBy("magnitude_bin") \
    .agg(count("*").alias("N")) \
    .withColumn("log_N", spark_log10(col("N")))

# Collect for numpy regression
data = mag_counts.orderBy("magnitude_bin").toPandas()
M = data["magnitude_bin"].values
logN = data["log_N"].values
b, a = np.polyfit(M, logN, 1)   # b-value (negative slope)
print(f"Gutenberg-Richter: a={a:.2f}, b={-b:.3f}")

# Exceedance probability for M >= 7.0 in 50 years
# P(exceedance) = 1 - exp(-lambda * T)
lambda_annual = 10 ** (a + b * 7.0)  # events per year
T = 50.0
prob_50yr = 1 - np.exp(-lambda_annual * T)
print(f"P(M>=7.0 in 50 yr) = {prob_50yr:.3f}")

# Uncertainty on b-value (Aki 1965 estimator)
N_events = int(catalog.filter(col("magnitude") >= 2.0).count())
M_mean   = float(catalog.filter(col("magnitude") >= 2.0)
                 .agg({"magnitude": "mean"}).collect()[0][0])
M_min    = 2.0
sigma_b  = b / np.sqrt(N_events)
print(f"b-value uncertainty (1-sigma): {sigma_b:.4f}")
\end{lstlisting}

%----------- section ----------->
\section{Databricks: Cloud-Native Geophysics}

\subsection{From Raw Waves to Real-Time Alerts: The Medallion Architecture}

In geophysics, the Medallion Architecture on Databricks streamlines the transition from raw seismic sensor feeds to predictive insights by organizing data into three functional stages, Figure~\ref{fig06}. The process begins with the \textbf{Bronze layer}, which ingests raw seismic waveforms from global APIs like USGS or local station streams, preserving the original signal. \textcolor{blue}{This is refined into the} \textbf{Silver layer}, where signal processing removes cultural noise \textcolor{blue}{(traffic, industrial vibrations)} to isolate tectonic activity. \textcolor{blue}{Finally,} the \textbf{Gold layer} provides curated, analytics-ready datasets used for magnitude and location prediction.

\begin{figure}[H]
    \centering
    \includegraphics[width=\textwidth]{Fig_06.png} 
    \caption{\textcolor{blue}{Medallion Architecture for} Geophysical Data Processing\textcolor{blue}{:} Istanbul EEW \& Seismic Analysis. Diagram source: authors.}
    \label{fig06}
\end{figure}

\textcolor{blue}{The key geophysical equations underpinning each layer are as follows.}

\textcolor{blue}{In the Silver layer, the signal-to-noise ratio (SNR) is computed over each window:}
\begin{equation}
\textcolor{blue}{\text{SNR}_{dB} = 10 \log_{10}\left(\frac{P_{\text{signal}}}{P_{\text{noise}}}\right)}
\end{equation}

\textcolor{blue}{where $P_{\text{signal}}$ and $P_{\text{noise}}$ are the signal and noise power estimated from pre- and post-event windows, respectively.}

\textcolor{blue}{The Brune source model, widely used in the Gold layer for spectral magnitude estimation, defines the far-field displacement spectrum as:}
\begin{equation}
\textcolor{blue}{\Omega(f) = \frac{\Omega_0}{1 + (f/f_c)^2}}
\end{equation}

\textcolor{blue}{where $\Omega_0$ is the long-period spectral level, and $f_c$ is the corner frequency related to the seismic moment $M_0$ and stress drop $\Delta\sigma$:}
\begin{equation}
\textcolor{blue}{f_c = 0.49 V_S \left(\frac{\Delta\sigma}{M_0}\right)^{1/3}}
\end{equation}

\textcolor{blue}{Moment magnitude $M_w$ is derived from the seismic moment as:}
\begin{equation}
\textcolor{blue}{M_w = \frac{2}{3}\log_{10}(M_0) - 10.7}
\end{equation}

\subsection{\textcolor{blue}{Geophysical Data} Hubs in Türkiye: Distributed Architecture}

The distributed approach to data processing corresponds to its current organisational structure, Table~\ref{tab03}:

\keepXColumns
\begin{tabularx}{\textwidth}{@{} 
    >{\hsize=0.6\hsize\raggedright\arraybackslash}X 
    >{\hsize=0.8\hsize\raggedright\arraybackslash}X 
    >{\hsize=1.6\hsize\raggedright\arraybackslash}X 
@{}}
    \caption{Principal Organizations and Big Data Activities in Marmara Region Seismology and EEW} \label{tab03} \\
    
    \toprule
    \textbf{Data Centre} & \textbf{Primary Focus} & \textbf{Seismic Big Data \& EEW Activities} \\ \midrule
    \endfirsthead
    
    \multicolumn{3}{c}{{\bfseries \tablename\ \thetable{} -- Continued from previous page}} \\
    \toprule
    \textbf{Data Centre} & \textbf{Primary Focus} & \textbf{Seismic Big Data \& EEW Activities} \\ \midrule
    \endhead
    
    \midrule
    \multicolumn{3}{r}{{Continued on next page...}} \\
    \endfoot
    
    \bottomrule
    \endlastfoot

    \textbf{KOERI (Kandilli)} \newline \textit{Est. 1868} & Operates the National Earthquake Monitoring Center (NEMC). & Processes real-time streams for the Istanbul EEW system, utilizing borehole and surface sensors to trigger automated gas/rail shutdowns. \\ \addlinespace
    
    \textbf{AFAD} \newline \textit{Est. 2009} & Manages the TDVMS with 1,100+ stations. & Acts as the official authority for high-velocity acceleration data used in real-time structural health monitoring. \\ \addlinespace
    
    \textbf{MTA} \newline \textit{Est. 1935 (Ankara)} & Curates the Active Fault Map of Türkiye. & Big data activities focus on GIS integration of paleoseismological data and surface rupture mapping. \\ \addlinespace
    
    \textbf{ITU Geophysics} \newline \textit{Est. 1952} & Houses the N. Canıtez Lab. & Pioneers AI-driven feature extraction for Marmara-specific site-response and tomography. \\ \addlinespace
    
    \textbf{METU Earth Systems} \newline \textit{Est. 2002} & Integrates multi-disciplinary big data. & Blends GNSS/GPS crustal deformation with seismic catalogs to quantify slip rate and strain on the NAFZ. \\ 

\end{tabularx}

\subsection{\textcolor{blue}{Governance,} Integration Planning\textcolor{blue}{, and Uncertainty Management}}

Governance and integration for earthquake monitoring must prioritize reliability, security, and interoperability by centralizing control through Unity Catalog. \textcolor{blue}{A key component is the systematic propagation of uncertainty from raw data through to final hazard products.} The implementation requires hardened security—using Private Links and Identity Federation—alongside Databricks Auto Loader for resilient ingestion of high-frequency station streams.

\textcolor{blue}{Confidence intervals for magnitude estimates are derived from the inter-station standard deviation $\sigma_M$:}
\begin{equation}
\textcolor{blue}{M \pm z_{\alpha/2} \cdot \frac{\sigma_M}{\sqrt{N_s}}}
\end{equation}
\textcolor{blue}{where $N_s$ is the number of contributing stations and $z_{\alpha/2}$ is the critical value for the desired confidence level (e.g., $z_{0.025} = 1.96$ for 95\% CI).}

\textcolor{blue}{Epicentral location uncertainty is quantified through the covariance matrix of the least-squares hypocenter solution:}
\begin{equation}
\textcolor{blue}{\mathbf{C}_{\mathbf{x}} = \sigma_t^2 (\mathbf{A}^T \mathbf{A})^{-1}}
\end{equation}
\textcolor{blue}{where $\mathbf{A}$ is the Jacobian matrix of arrival time partial derivatives with respect to hypocentral coordinates, and $\sigma_t$ is the arrival time picking uncertainty.}

\section{Deterministic Infrastructure and Performance Optimization for EEW Systems}

\subsection{Computational Tiering and Latency Requirements}

EEW systems rely on a tiered computational architecture with increasingly intensive processing demands and stringent latency constraints. The primary workload involves continuous, sub-second ingestion and filtering of high-velocity seismic sensor telemetry, requiring deterministic I/O throughput. \textcolor{blue}{This is followed by event detection and phase picking, where ML models must scale instantly to identify P-wave arrivals across thousands of concurrent streams.}

\textcolor{blue}{The minimum warning time $T_w$ at distance $r$ from a fault is bounded by:}
\begin{equation}
\textcolor{blue}{T_w = \frac{r}{V_S} - T_{\text{proc}} - T_{\text{dissem}}}
\end{equation}
\textcolor{blue}{where $T_{\text{dissem}}$ is the alert dissemination time (typically 1--2~s for cellular broadcast). For the Kadıköy district of Istanbul ($r \approx 55$~km, $V_S = 3.5$~km/s), $T_w \approx 13.7$~s assuming $T_{\text{proc}} = 2$~s.}

\subsection{GPU-Accelerated Systems for Marmara Region Earthquake Alerts}

To optimize seismic big data for the Istanbul and Marmara region, the strategy focuses on a tiered approach balancing high-speed performance with cost sustainability. By implementing edge-based signal processing at seafloor stations, the system filters raw noise at the source, reducing data transfer bottlenecks. During seismic swarms, the architecture utilizes programmatic JSON configurations to trigger GPU-accelerated pools, preventing memory spills and ensuring sub-second phase picking, Figure~\ref{fig08}.

\begin{figure}[H]
    \begin{center}
	\hspace{-35pt}\includegraphics[width=13.0cm]{Fig_08.png}
    \end{center}
    \caption{Earthquake Early Warning (EEW) system data flow on Databricks. Diagram source: authors.}
    \label{fig08}
\end{figure}

\textcolor{blue}{The performance of the EEW framework is evaluated against three key metrics. The Unit Cost of Detection (UCD) is defined as:}
\begin{equation}
\textcolor{blue}{\text{UCD} = \frac{C_{\text{compute}} + C_{\text{storage}} + C_{\text{egress}}}{N_{\text{detected}}} \quad [\$/\text{event}]}
\end{equation}

\textcolor{blue}{Minimizing UCD while maintaining a recall $>$~0.99 for $M \geq 3.0$ events constitutes the primary optimization objective for the Marmara Sea EEW.}

\subsection{Hybrid Cluster Architecture for Sub-Second Earthquake Detection}

In the Marmara Sea region, an EEW-optimized Databricks cluster \textcolor{blue}{operates as} a mission-critical engine within a hybrid mosaic architecture. Azure Databricks serves as a managed, high-performance platform built upon the open-source Apache Spark engine. Spark provides the core distributed computing power to parallelize seafloor telemetry across a cluster, while Databricks enhances this with a managed runtime that optimizes memory management and I/O throughput.

By utilizing programmatic JSON configurations and GPU-accelerated pools, the system handles seismic swarms with sub-second latency, preventing fatal memory spills during the ``Golden Seconds''. The workflow integrates seafloor smart nodes for edge-based ingestion (Bronze), Physics-Informed Neural Networks for real-time denoising (Silver), and refined Digital Twin models for hazard mitigation (Gold).

%----------- section ----------->
\section{Conclusions}

\textcolor{blue}{This review has examined the application of big data architectures and AI methodologies to seismic hazard assessment in the Marmara region of Türkiye.} \textcolor{blue}{Several principal findings emerge.} \textcolor{blue}{First, we established that EEW and earthquake forecasting are fundamentally distinct problems requiring separate data pipelines, latency budgets, and validation frameworks.} \textcolor{blue}{EEW demands deterministic sub-second processing, whereas probabilistic forecasting requires decadal seismic catalogs and uncertainty quantification.}

\textcolor{blue}{Second, the Gutenberg--Richter analysis of the AFAD catalog (Script~8) yields a b-value consistent with the regional NAFZ seismicity, and the 50-year exceedance probability for $M \geq 7.0$ events underscores the urgency of operational EEW in Istanbul. The statistical validation framework (Equations~1--4, Scripts~4 and~8) provides quantitative metrics—RMSE, precision, recall, and confidence intervals—that transform the analysis from descriptive to rigorously scientific.}

\textcolor{blue}{Third,} the successful implementation of cloud optimization for the Istanbul EEW requires a multi-layered communication strategy \textcolor{blue}{that balances} technical \textcolor{blue}{performance} and fiscal \textcolor{blue}{sustainability}. By \textcolor{blue}{adopting} FinOps-driven fiscal predictability\textcolor{blue}{,} ``Golden Path'' technical agility\textcolor{blue}{,} and automated Policy-as-Code (PaC) compliance\textcolor{blue}{,} the organization \textcolor{blue}{can resolve the tension} between rapid innovation and strict KVKK data sovereignty.

\textcolor{blue}{Fourth, the eight Spark-based code scripts presented here illustrate the complete pipeline from raw MiniSEED ingestion through Butterworth filtering, STA/LTA event detection, PhaseNet phase picking, magnitude estimation with statistical validation, ambient noise cross-correlation, and real-time alert latency monitoring. Together, these components constitute a reproducible and scalable framework for Marmara Sea seismic monitoring.}

\textcolor{blue}{The primary limitation of this review is the absence of a fully deployed, end-to-end operational test using live AFAD/KOERI streams; such validation constitutes the immediate next step. Future work should also address the systematic characterization of anthropogenic noise sources in the Istanbul metropolitan area—including metro, ferry, and industrial signals—and their impact on EEW false alarm rates. Additionally, integrating GNSS geodetic data for real-time strain monitoring alongside seismic streams represents a high-priority direction for improving probabilistic forecasting in the Marmara region.}

%----------- section ----------->
\authorcontributions{Supervision, conceptualization, administration, funding acquisition, and project administration -- A.C.Z. Methodology, software, resources, writing---original draft preparation, data curation, visualization, formal analysis, validation, writing---review and editing, and investigation, P.L.}

\institutionalreview{Not applicable.}

\informedconsent{Not applicable.}

\dataavailability{Not applicable.} 

\acknowledgments{The authors thank the reviewers for their careful reading and constructive comments on this manuscript.}

\conflictsofinterest{The authors declare no conflict of interest.} 

\abbreviations{Abbreviations}{
The following abbreviations are used in this manuscript:\\

\noindent 
\begin{tabular}{@{}ll}
AFAD & Disaster and Emergency Management Authority \\
AI & Artificial Intelligence \\
API & Application Programming Interface \\
CCS & Carbon Capture and Storage \\
CNN & Convolutional Neural Networks\\
DAS & Distributed Acoustic Sensing \\
DAG & Directed Acyclic Graph \\
DL & Deep Learning \\
EEW & Earthquake Early Warning \\
ETL & Extract, Transform, and Load \\
FWI & Full Waveform Inversion \\
GPU & Graphics Processing Unit \\
GNNs & Graph Neural Networks \\
HPC & High-Performance Computing \\
I/O & Input-Output \\
KOERI & Kandilli Observatory and Earthquake Research Institute\\
KVKK & Kişisel Verileri Korunması Kanunu \\
LSTM & Long Short-Term Memory \\
ML & Machine Learning \\
NAF & North Anatolian Fault \\
NAFZ & North Anatolian Fault Zone \\
NEMC & National Earthquake Monitoring Center \\
PINN & Physics-Informed Neural Network \\
PaC & Policy-as-Code \\
QC & Quality Control \\
RMSE & Root Mean Square Error \\
RNN & Recurrent Neural Network \\
SaaS & Software as a Service \\
SEG-Y & Society of Exploration Geophysicists -- Y format \\
SHRA & Seismic Hazard Risk Assessment \\
SNR & Signal-to-Noise Ratio \\
SQL & Structured Query Language \\
STA/LTA & Short-Term Average / Long-Term Average \\
UCD & Unit Cost of Detection \\
\end{tabular}
}

%%%%%%%%%%%%%%%%%%%%%%%%%%%%%%%%%%%%%%%%%%
\begin{adjustwidth}{-\extralength}{0cm}

\reftitle{References}
\nocite{*}

\bibliography{Bibliography.bib}

%%%%%%%%%%%%%%%%%%%%%%%%%%%%%%%%%%%%%%%%%%
\end{adjustwidth}
\end{document}
